% =====================================================================
%  PRESENTATION TIPE (Beamer) - Filiere MPSI/MP
%  Modele poutre 1D. .
%
%  Compilation : pdflatex presentation_tipe.tex   (x2 pour la nav.)
%  Les figures fig_poutre_*.png sont produites par simulation_poutre.py.
%  Un cadre de substitution s'affiche tant qu'elles n'existent pas.
% =====================================================================
\documentclass[11pt,aspectratio=169]{beamer}
\usepackage{tikz}
\usetikzlibrary{arrows.meta,positioning,patterns,shapes.geometric}
\definecolor{verre}{RGB}{30,120,180}
\definecolor{danger}{RGB}{200,40,40}
\definecolor{ok}{RGB}{30,150,90}

\usetheme{Madrid}
\usecolortheme{whale}
\setbeamertemplate{navigation symbols}{}
\setbeamertemplate{caption}[numbered]

\usepackage[utf8]{inputenc}
\usepackage[T1]{fontenc}
\usepackage[french]{babel}
\usepackage{amsmath,amssymb}
\usepackage{graphicx}
\usepackage{tikz}
\usetikzlibrary{arrows.meta,positioning,decorations.pathmorphing,
                patterns,shapes.geometric,calc,backgrounds}

% ---- couleurs ----
\definecolor{verre}{RGB}{30,120,180}
\definecolor{danger}{RGB}{200,40,40}
\definecolor{ok}{RGB}{30,150,90}
\definecolor{gristik}{RGB}{120,120,120}

% ---- raccourcis TikZ reutilisables ----
% sol hachure
\newcommand{\sol}[2]{% #1 xgauche  #2 xdroite
  \draw[thick](#1,0)--(#2,0);
  \foreach \xx in {#1,...,#2}{\draw[gristik](\xx,0)--(\xx-0.25,-0.25);}}
% ressort horizontal entre (a) et (b)
% smartphone (rectangle arrondi) centre en (#1) de demi-largeur 0.35 hauteur 0.7
\newcommand{\phone}[3]{% #1 x  #2 y  #3 couleur
  \draw[#3,fill=#3!12,thick,rounded corners=2pt]
       (#1-0.32,#2-0.62) rectangle (#1+0.32,#2+0.62);
  \draw[#3!60] (#1-0.24,#2-0.45) rectangle (#1+0.24,#2+0.45);}

% inclusion robuste d'une figure de simulation
\newcommand{\figauto}[3]{%
  \IfFileExists{#1}{\includegraphics[width=#2\linewidth]{#1}}{%
    \fcolorbox{gristik}{gray!8}{\parbox[c][3cm][c]{#2\linewidth}{\centering\footnotesize
    \textit{[\texttt{#1} : lancer \texttt{simulation\_poutre.py}]}\\[2pt]#3}}}}

% ---- page de titre ----
\title[Écran en verre trempé]{Optimiser l'épaisseur d'un écran protecteur\\en verre trempé pour qu'il casse moins}
\subtitle{Modèle « poutre »}
\author{TIPE}
\institute{Filière MPSI/MP}
\date{2026-2027}


\begin{document}



% ===================== Slide 1 : titre =====================
{\setbeamertemplate{footline}{}
\begin{frame}[plain]
  \titlepage
  \vspace{-0.6cm}
  \begin{center}
  \begin{tikzpicture}[scale=0.9]
    % telephone qui tombe vers un sol, ecran fissure
    \sol{-1}{4}
    \phone{1.5}{1.3}{verre}
    \draw[-{Latex[length=3mm]},danger,very thick](1.5,2.4)--(1.5,2.0);
    \node[danger] at (3.1,2.2){\footnotesize chute};
    % petite fissure au coin
    \draw[danger,thick](1.18,0.72)--(1.32,0.95)--(1.22,1.15)--(1.36,1.35);
  \end{tikzpicture}
  \end{center}
\end{frame}
}

% ===================== Slide 2 : le constat =====================
\begin{frame}{Le problème : ça casse,presque toujours au même endroit}
  \begin{columns}[T]
  \column{0.52\textwidth}
    \begin{itemize}
      \item Un smartphone tombe en moyenne plusieurs fois par an.
      \item Les écrans en verre trempé sont durs\dots\ mais cassants.
      \item Les fissures partent presque toujours d'un \textbf{coin} ou d'un \textbf{bord}.
    \end{itemize}
    \vspace{0.3cm}
    \begin{block}{Intuition de départ}
      Et si on mettait \textbf{plus de verre là où ça casse} (les bords)
      et \textbf{moins au centre} (pour garder le toucher) ?
    \end{block}
  \column{0.46\textwidth}
    \centering
    \begin{tikzpicture}[scale=1.05]
      \draw[verre,thick,fill=verre!8,rounded corners=4pt](0,0) rectangle (3,5);
      % zone tactile centre
      \draw[ok,dashed](1.5,2.5) circle (1);
      \node[ok] at (1.5,2.5){\scriptsize zone tactile};
      % fissures aux coins
      \foreach \x/\y in {0.1/0.1, 2.9/0.1, 0.1/4.9, 2.9/4.9}{
        \draw[danger,thick](\x,\y)
          .. controls ($(\x,\y)+(0.4,0.3)$) .. ($(\x,\y)+(0.7,0.7)$);}
      \node[danger] at (1.5,-0.45){\scriptsize fissures aux coins};
    \end{tikzpicture}
  \end{columns}
\end{frame}

% ===================== Slide 3 : problematique =====================
\begin{frame}{La question qui se pose:}
  \begin{center}
  \begin{beamercolorbox}[wd=0.92\textwidth,sep=10pt,rounded=true,shadow=true]{block body}
  \centering\large
  \textbf{Quel profil d'épaisseur de l'écran rend la rupture\\
  la moins probable lors d'une chute,}\\[4pt]
  \normalsize tout en gardant un centre fin ($\le 0{,}5$~mm) pour le toucher
  et une masse raisonnable ?
  \end{beamercolorbox}
  \end{center}
  \vspace{0.4cm}
  \begin{center}
  \begin{tikzpicture}[scale=1]
    % trois coupes d'ecran : uniforme / parabolique / expo
    \begin{scope}
      \draw[verre,fill=verre!15](0,0) rectangle (2.4,0.4);
      \node at (1.2,-0.35){\scriptsize uniforme ?};
    \end{scope}
    \begin{scope}[xshift=3.3cm]
      \fill[verre!15] (0,0) .. controls (1.2,0.55) .. (2.4,0)
                       -- (2.4,0.18) .. controls (1.2,0.6) .. (0,0.18) -- cycle;
      \draw[verre] (0,0.18) .. controls (1.2,0.6) .. (2.4,0.18);
      \draw[verre] (0,0)--(2.4,0);
      \node at (1.2,-0.35){\scriptsize  plus épais au centre ?};
    \end{scope}
    \begin{scope}[xshift=6.6cm]
      \fill[verre!15] (0,0) rectangle (2.4,0.2);
      \fill[verre!15] (0,0.2) .. controls (1.2,0.0) .. (2.4,0.2) -- (2.4,0.2)--(0,0.2);
      \draw[verre](0,0.6)--(0,0.2) (2.4,0.6)--(2.4,0.2);
      \draw[verre](0,0.6) .. controls (1.2,0.2) .. (2.4,0.6);
      \draw[verre](0,0)--(2.4,0);
      \node at (1.2,-0.35){\scriptsize  plus épais aux bords ?};
    \end{scope}
  \end{tikzpicture}
  \end{center}
\end{frame}

% ===================== Slide 4 : la demarche =====================
\begin{frame}{Notre démarche en une image}
  \begin{center}
  \begin{tikzpicture}[node distance=0.5cm and 0.7cm,
      box/.style={draw,rounded corners,fill=verre!10,align=center,
                  minimum height=1.1cm,text width=2.5cm,font=\footnotesize},
      ar/.style={-{Latex},thick,verre}]
    \node[box](m){\textbf{1. Modèle}\\ poutre + énergie\\ $\sigma \propto 1/\sqrt{e}$};
    \node[box,right=of m](v){\textbf{2. Validation}\\ analogie circuit RLC};
    \node[box,right=of v](s){\textbf{3. Simulation}\\ Monte-Carlo\\ $P(\text{rupture})$};
    \node[box,right=of s](e){\textbf{4. Expérience}\\ son + FFT\\ carte de rigidité};
    \draw[ar](m)--(v); \draw[ar](v)--(s); \draw[ar](s)--(e);
  \end{tikzpicture}
  \end{center}
  \vspace{0.3cm}
  \centering\small Un seul fil conducteur : \textcolor{verre}{comprendre} $\rightarrow$
  \textcolor{verre}{vérifier} $\rightarrow$ \textcolor{verre}{quantifier} $\rightarrow$
  \textcolor{verre}{tester}.
\end{frame}


% ===================== Slide 5 : la chute =====================
\begin{frame}{Étape 1 --- La chute donne de l'énergie}
  \begin{columns}[c]
  \column{0.5\textwidth}
    \begin{tikzpicture}[scale=1]
      \draw[-{Latex}](0,5)--(0,-0.2) node[right]{};
      \node at (-0.35,2.5){\rotatebox{90}{\footnotesize hauteur $h$}};
      \phone{1.6}{4.2}{verre}
      \draw[-{Latex[length=2.5mm]},danger,thick](1.6,3.4)--(1.6,2.6)
            node[midway,right]{\footnotesize $v$};
      \phone{1.6}{0.7}{gristik}
      \sol{0.4}{3}
      \draw[{Latex}-{Latex}](0.15,0)--(0.15,4.2) node[midway,left]{};
    \end{tikzpicture}
  \column{0.5\textwidth}
    \begin{block}{Mécanique du point (au programme)}
    Chute libre d'une hauteur $h$ :
    \[ v=\sqrt{2gh} \]
    Énergie cinétique juste avant l'impact :
    \[ E_c=\tfrac12\,m\,v^{2}=m\,g\,h \]
    \end{block}
    \footnotesize\textbf{Exemple :} $h=1$~m $\Rightarrow v\approx 4{,}4$~m/s,
    $E_c\approx 1{,}8$~J pour $m=180$~g.\\[4pt]
    \textcolor{danger}{Toute cette énergie doit être « encaissée » en un instant.}
  \end{columns}
\end{frame}

% ===================== Slide 6 : le verre = un ressort =====================
\begin{frame}{Étape 2 --- Le verre se comporte comme un ressort}
  \begin{columns}[c]
  \column{0.46\textwidth}
    \begin{tikzpicture}[scale=1]
      % mur
      \draw[thick](0,-0.2)--(0,2.2);
      \foreach \yy in {-0.1,0.2,...,2.1}{\draw[gristik](0,\yy)--(-0.25,\yy-0.2);}
      % ressort
      \draw[decorate,decoration={coil,aspect=0.6,segment length=4pt,amplitude=4pt},verre,thick]
            (0,1)--(2.2,1);
      % masse
      \draw[fill=verre!15,thick](2.2,0.4) rectangle (3.2,1.6);
      \node at (2.7,1){\footnotesize $m$};
      \draw[-{Latex},danger,thick](3.3,1)--(4.0,1) node[right]{\footnotesize $x$};
    \end{tikzpicture}
  \column{0.54\textwidth}
    \begin{itemize}
      \item Sous le choc, l'écran fléchit puis revient : c'est un
            \textbf{oscillateur}.
      \item Il stocke une énergie élastique :
      \[ E_{\text{stockée}}=\tfrac12\,k\,x^{2} \]
      \item $k$ = \textbf{raideur} de l'écran (sa « dureté à plier »),
            $x$ = fléche (de combien il s'enfonce).
    \end{itemize}
    \vspace{0.2cm}
    \centering\footnotesize\textcolor{verre}{Tout le problème : relier $k$ et $x$
    à l'épaisseur $e$.}
  \end{columns}
\end{frame}

% ===================== Slide 7 : modele poutre 1D =====================
\begin{frame}{Étape 3 --- On modélise l'écran par une \emph{poutre}}
  \begin{columns}[c]
  \column{0.54\textwidth}
    \begin{tikzpicture}[scale=1]
      % appuis
      \fill[gristik](0,0)--(-0.3,-0.4)--(0.3,-0.4)--cycle;
      \fill[gristik](5,0)--(4.7,-0.4)--(5.3,-0.4)--cycle;
      % poutre droite (pointilles) et flechie
      \draw[gristik,dashed](0,0)--(5,0);
      \draw[verre,very thick](0,0) .. controls (2.5,-1.1) .. (5,0);
      % charge
      \draw[-{Latex[length=2.5mm]},danger,thick](2.5,1)--(2.5,-0.5)
            node[midway,right]{\footnotesize $F$ (choc)};
      % fleche
      \draw[{Latex}-{Latex}](5.4,0)--(5.4,-0.85) node[midway,right]{\footnotesize $x$};
      % portee
      \draw[{Latex}-{Latex}](0,-0.7)--(5,-0.7) node[midway,below]{\footnotesize portée $L$};
    \end{tikzpicture}
  \column{0.46\textwidth}
    \begin{block}{Simplification clé (1D)}
    On remplace la plaque (2D) par une \textbf{poutre sur deux appuis}
    qui se courbe dans \emph{une} seule direction.
    \end{block}
    \footnotesize
    \begin{itemize}
      \item Beaucoup plus simple : c'est de la \textbf{RDM de base}.
      \item On garde l'essentiel : comment $k$ dépend de $e$.
      \item On surestime un peu $\sigma$ $\Rightarrow$ choix
            \textcolor{ok}{prudent}.
    \end{itemize}
  \end{columns}
\end{frame}

% ===================== Slide 8 : raideur ~ e^3 =====================
\begin{frame}{Étape 4 --- Plus c'est épais, beaucoup plus c'est rigide}
  \begin{columns}[c]
  \column{0.4\textwidth}
    \centering
    \begin{tikzpicture}[scale=1]
      % section rectangulaire b x e
      \draw[verre,thick,fill=verre!12](0,0) rectangle (2.2,0.7);
      \draw[{Latex}-{Latex}](0,-0.25)--(2.2,-0.25) node[midway,below]{\footnotesize largeur $b$};
      \draw[{Latex}-{Latex}](2.45,0)--(2.45,0.7) node[midway,right]{\footnotesize épaisseur $e$};
      \node at (1.1,1.2){\footnotesize section de la poutre};
    \end{tikzpicture}
  \column{0.6\textwidth}
    \begin{block}{RDM : moment quadratique et raideur}
    \[ I=\frac{b\,e^{3}}{12} \qquad
       k=\frac{48\,E\,I}{L^{3}}=\frac{4\,E\,b}{L^{3}}\;e^{3} \]
    \end{block}
    \begin{center}
    \fbox{\textcolor{verre}{$\;k \propto e^{3}\;$}}\quad
    \footnotesize doubler l'épaisseur $\Rightarrow$ raideur $\times 8$ !
    \end{center}
    \footnotesize $E$ : rigidité du verre (module d'Young, $\approx 73$~GPa).
  \end{columns}
\end{frame}

% ===================== Slide 9 : energie -> fleche =====================
\begin{frame}{Étape 5 --- De l'énergie à la fléche}
  \begin{columns}[c]
  \column{0.5\textwidth}
    \begin{tikzpicture}[scale=0.95]
      % barre energie chute
      \draw[fill=danger!25](0,0) rectangle (1,3);
      \node[align=center] at (0.5,-0.4){\scriptsize énergie\\\scriptsize de chute $E_c$};
      \draw[-{Latex},thick](1.3,1.5)--(2.6,1.5) node[midway,above]{\scriptsize fraction $\eta$};
      % barre energie flexion
      \draw[fill=verre!30](2.9,0) rectangle (3.9,1.0);
      \node[align=center] at (3.4,-0.4){\scriptsize énergie\\\scriptsize de flexion};
      % reste
      \node[align=left,gristik] at (5.6,1.8){\scriptsize le reste :\\\scriptsize rebond, son,\\\scriptsize corps du tél.};
    \end{tikzpicture}
  \column{0.5\textwidth}
    \begin{block}{Conservation de l'énergie}
    Une fraction $\eta$ de l'énergie de chute sert à plier le verre :
    \[ \tfrac12\,k\,x_{\max}^{2}=\eta\,E_c
       \;\Rightarrow\; x_{\max}=\sqrt{\dfrac{2\,\eta\,E_c}{k}} \]
    \end{block}
    \footnotesize
    \textcolor{danger}{\textbf{Hypothèse à défendre :}} on \emph{admet}
    $\eta\approx 5\%$ (le reste part en rebond, son, déformation du
    téléphone). Cela évite le calcul compliqué du contact.
  \end{columns}
\end{frame}

% ===================== Slide 10 : fleche -> contrainte =====================
\begin{frame}{Étape 6 --- De la fléche à la contrainte qui casse}
  \begin{columns}[c]
  \column{0.5\textwidth}
    \begin{tikzpicture}[scale=1]
      \draw[verre,very thick](0,0) .. controls (2,-1) .. (4,0);
      % face du haut comprimee, face du bas tendue
      \node[ok] at (2,0.35){\footnotesize compression (dessus)};
      \draw[-{Latex},danger](0.6,-0.75)--(0.0,-0.95);
      \draw[-{Latex},danger](3.4,-0.75)--(4.0,-0.95);
      \node[danger] at (2,-1.15){\footnotesize traction (dessous) $\rightarrow$ c'est ici que ça casse};
    \end{tikzpicture}
  \column{0.5\textwidth}
    \begin{block}{Contrainte de flexion (RDM)}
    \[ \sigma_{\max}=\frac{6\,E\,e}{L^{2}}\;x_{\max} \]
    \end{block}
    \footnotesize Le verre casse en \textbf{traction}, sur la face opposée
    au choc. Plus la fléche $x_{\max}$ est grande, plus $\sigma_{\max}$ est
    grande.
  \end{columns}
\end{frame}

% ===================== Slide 11 : resultat cle sigma ~ 1/sqrt(e) =====================
\begin{frame}{Le résultat central : $\boxed{\sigma_{\max}\propto 1/\sqrt{e}}$}
  \begin{columns}[c]
  \column{0.5\textwidth}
    En combinant les étapes 4, 5 et 6 :
    \[ \sigma_{\max}=\frac{6Ee}{L^2}\sqrt{\frac{2\eta E_c}{k}},
       \quad k\propto e^{3} \]
    \[ \Rightarrow\quad \sigma_{\max}\;\propto\;\frac{1}{\sqrt{e}} \]
    \vspace{0.2cm}
    \begin{beamercolorbox}[sep=6pt,rounded=true]{block body}
    \centering\textbf{Épaissir localement $\Rightarrow$ moins de contrainte
    $\Rightarrow$ moins de casse.}
    \end{beamercolorbox}
  \column{0.5\textwidth}
    \centering
    \begin{tikzpicture}[scale=1]
      \draw[-{Latex}](0,0)--(4.2,0) node[right]{\footnotesize $e$};
      \draw[-{Latex}](0,0)--(0,3.2) node[above]{\footnotesize $\sigma_{\max}$};
      \draw[verre,very thick,domain=0.4:4,samples=60,smooth]
            plot (\x,{1.6/sqrt(\x)});
      \draw[danger,dashed](0,1.1)--(4,1.1) node[right]{\footnotesize seuil};
      \node[verre] at (2.6,1.9){\footnotesize $1/\sqrt{e}$};
    \end{tikzpicture}
    \\[2pt]\footnotesize même conclusion que le modèle « plaque », en bien plus simple.
  \end{columns}
\end{frame}

% ===================== Slide 12 : critere de rupture =====================
\begin{frame}{Quand est-ce que ça casse ?}
  \begin{columns}[c]
  \column{0.5\textwidth}
    \begin{block}{Critère simple}
    Le verre casse si la contrainte dépasse un \textbf{seuil} $\sigma_c$ :
    \[ \sigma_{\max} > \sigma_c \;\Rightarrow\; \text{rupture} \]
    \end{block}
    \footnotesize
    \begin{itemize}
      \item Le verre est \textbf{fragile} : pas de déformation prévenante.
      \item Le seuil $\sigma_c$ n'est pas fixe : il \textbf{varie} d'un point
            à l'autre (micro-défauts). On le modélise par une loi
            \textbf{gaussienne} (au programme) autour de $\approx 850$~MPa.
    \end{itemize}
  \column{0.5\textwidth}
    \centering
    \begin{tikzpicture}[scale=1]
      \draw[-{Latex}](-0.2,0)--(4.4,0) node[right]{\footnotesize $\sigma$};
      \draw[-{Latex}](0,0)--(0,2.6);
      % gaussienne du seuil
      \draw[verre,thick,domain=0.4:4,samples=80,smooth]
        plot (\x,{2*exp(-(\x-2.4)*(\x-2.4)/0.35)});
      \node[verre] at (2.4,2.25){\footnotesize seuil $\sigma_c$};
      % contrainte appliquee
      \draw[danger,very thick](1.5,0)--(1.5,2.2) node[above]{\footnotesize $\sigma_{\max}$};
      \fill[danger!20,domain=0.4:1.5,samples=40] plot (\x,{2*exp(-(\x-2.4)*(\x-2.4)/0.35)})
            -- (1.5,0) -- (0.4,0) -- cycle;
      \node[danger,align=center] at (0.95,-0.55){\scriptsize P(casse)};
    \end{tikzpicture}
  \end{columns}
\end{frame}


% ===================== Slide 13 : les 3 profils =====================
\begin{frame}{Les trois profils que l'on compare}
  % --- rangee 1 : les trois profils, chacun dans sa colonne ---
  \begin{columns}[T]
    \column{0.33\textwidth}
      \centering
      \textbf{\footnotesize Uniforme}\\[4pt]
      \begin{tikzpicture}[scale=0.8]
        \draw[verre,thick,fill=verre!12](0,0) rectangle (2.6,0.4);
      \end{tikzpicture}\\[4pt]
      {\scriptsize $e(s)=e_0$}
    \column{0.33\textwidth}
      \centering
      \textbf{\footnotesize Parabolique}\\[4pt]
      \begin{tikzpicture}[scale=0.8]
        \fill[verre!12](0,0)--(2.6,0)--(2.6,0.7)
          .. controls (1.3,0.12) .. (0,0.7)--cycle;
        \draw[verre,thick](0,0.7) .. controls (1.3,0.12) .. (2.6,0.7);
        \draw[verre,thick](0,0)--(2.6,0)--(2.6,0.7) (0,0)--(0,0.7);
      \end{tikzpicture}\\[4pt]
      {\scriptsize $e(s)=e_c+(e_b-e_c)\,s^2$}
    \column{0.33\textwidth}
      \centering
      \textbf{\footnotesize Exponentiel}\\[4pt]
      \begin{tikzpicture}[scale=0.8]
        \fill[verre!12](0,0)--(2.6,0)--(2.6,0.85)
          .. controls (1.8,0.2) and (0.9,0.12) .. (0,0.28)--cycle;
        \draw[verre,thick](0,0.28) .. controls (0.9,0.12) and (1.8,0.2) .. (2.6,0.85);
        \draw[verre,thick](0,0)--(2.6,0)--(2.6,0.85) (0,0)--(0,0.28);
      \end{tikzpicture}\\[4pt]
      {\scriptsize $e(s)=e_c\,e^{\alpha s}$}
  \end{columns}
  \vspace{0.35cm}
  % --- rangee 2 : vue de dessus + explication de la variable s ---
  \begin{columns}[c]
    \column{0.28\textwidth}
      \centering
      \begin{tikzpicture}[scale=0.4]
        \fill[verre!30,even odd rule,rounded corners=5pt]
              (0,0) rectangle (3,6) (0.55,0.55) rectangle (2.45,5.45);
        \draw[verre,thick,rounded corners=5pt](0,0) rectangle (3,6);
        \draw[verre,thick,dashed,rounded corners=3pt](0.55,0.55) rectangle (2.45,5.45);
        
ode[font=\tiny,verre] at (1.5,3){$s=0$};
      \end{tikzpicture}\\[2pt]
      {\tiny vue de dessus : $s=1$ sur tout le bord}
    \column{0.72\textwidth}
      \footnotesize
      \textbf{On prend s le paramètre représentant la proximité au bords}
      \[ s=\max\!\left(\frac{|x|}{L/2},\ \frac{|y|}{\ell/2}\right)\in[0,1] \]
      $s=0$ au centre, $s=1$ sur \textbf{tout} le pourtour $\Rightarrow$
      surépaisseur \textbf{uniforme} le long des quatre bords.\\[2pt]
      {\scriptsize Contrainte : centre fin ($e_c\le 0{,}5$~mm),
      masse $\le 2{,}5\times$ l'uniforme.}
  \end{columns}
\end{frame}

% ===================== Slide (NOUVELLE) : perspective + centre, exponentiel corrige =====================
\begin{frame}{Les trois profils vus en perspective }
  \begin{columns}[T]

  % ---------- 1) UNIFORME ----------
  \column{0.33\textwidth}
    \centering \textbf{\footnotesize Uniforme}\\[6pt]
    \begin{tikzpicture}[scale=0.5]
      \fill[verre!18](0,0.5)--(4,0.5)--(4.9,1.05)--(0.9,1.05)--cycle;
      \fill[verre!32](4,0)--(4,0.5)--(4.9,1.05)--(4.9,0.55)--cycle;
      \fill[verre!12](0,0)--(4,0)--(4,0.5)--(0,0.5)--cycle;
      \draw[verre,thick](0,0)--(4,0)--(4,0.5)--(0,0.5)--cycle;
      \draw[verre,thick](4,0)--(4.9,0.55) (4,0.5)--(4.9,1.05) (0,0.5)--(0.9,1.05);
      \draw[verre,thick](0.9,1.05)--(4.9,1.05)--(4.9,0.55);
      \draw[danger,dashed,thick](2,0.5)--(2.9,1.05);
      \fill[danger](2.45,0.775) circle (0.09);
      
ode[danger,font=\tiny,above right] at (2.45,0.775){centre};
    \end{tikzpicture}\\[6pt]
    {\scriptsize épaisseur constante}

  % ---------- 2) PARABOLIQUE ----------
  \column{0.33\textwidth}
    \centering \textbf{\footnotesize Parabolique}\\[6pt]
    \begin{tikzpicture}[scale=0.5]
      \fill[verre!18](0,1.0) .. controls (1,0.45) and (3,0.45) .. (4,1.0)
            -- (4.9,1.55) .. controls (3.9,1.0) and (1.9,1.0) .. (0.9,1.55) -- cycle;
      \fill[verre!32](4,0)--(4,1.0)--(4.9,1.55)--(4.9,0.55)--cycle;
      \fill[verre!12](0,0)--(4,0)--(4,1.0)
            .. controls (3,0.45) and (1,0.45) .. (0,1.0)--cycle;
      \draw[verre,thick](0,0)--(4,0)--(4,1.0)
            .. controls (3,0.45) and (1,0.45) .. (0,1.0)--cycle;
      \draw[verre,thick](4,0)--(4.9,0.55)--(4.9,1.55);
      \draw[verre,thick](4,1.0)--(4.9,1.55) (0,1.0)--(0.9,1.55);
      \draw[verre,thick](0.9,1.55) .. controls (1.9,1.0) and (3.9,1.0) .. (4.9,1.55);
      \draw[danger,dashed,thick](2,0.5875)--(2.9,1.1375);
      \fill[danger](2.45,0.8625) circle (0.09);
      
ode[danger,font=\tiny,above right] at (2.45,0.8625){centre (fin)};
    \end{tikzpicture}\\[6pt]
    {\scriptsize épais aux deux bords, fin au centre}

  % ---------- 3) EXPONENTIEL (corrige : symetrique) ----------
  \column{0.33\textwidth}
    \centering \textbf{\footnotesize Exponentiel}\\[6pt]
    \begin{tikzpicture}[scale=0.5]
      % dessus (cuvette symetrique : fond plat, parois raides)
      \fill[verre!18]
        (0,1.0) .. controls (0.8,0.3) and (1.7,0.32) .. (2,0.32)
                .. controls (2.3,0.32) and (3.2,0.3) .. (4,1.0)
        -- (4.9,1.55)
        .. controls (4.1,0.85) and (3.2,0.87) .. (2.9,0.87)
        .. controls (2.6,0.87) and (1.7,0.85) .. (0.9,1.55) -- cycle;
      \fill[verre!32](4,0)--(4,1.0)--(4.9,1.55)--(4.9,0.55)--cycle;
      \fill[verre!12]
        (0,0)--(4,0)--(4,1.0)
        .. controls (3.2,0.3) and (2.3,0.32) .. (2,0.32)
        .. controls (1.7,0.32) and (0.8,0.3) .. (0,1.0)--cycle;
      \draw[verre,thick]
        (0,0)--(4,0)--(4,1.0)
        .. controls (3.2,0.3) and (2.3,0.32) .. (2,0.32)
        .. controls (1.7,0.32) and (0.8,0.3) .. (0,1.0)--cycle;
      \draw[verre,thick](4,0)--(4.9,0.55)--(4.9,1.55);
      \draw[verre,thick](4,1.0)--(4.9,1.55) (0,1.0)--(0.9,1.55);
      \draw[verre,thick]
        (0.9,1.55) .. controls (1.7,0.85) and (2.6,0.87) .. (2.9,0.87)
                   .. controls (3.2,0.87) and (4.1,0.85) .. (4.9,1.55);
      \draw[danger,dashed,thick](2,0.32)--(2.9,0.87);
      \fill[danger](2.45,0.595) circle (0.09);
      
ode[danger,font=\tiny,above right] at (2.45,0.595){centre (fin)};
    \end{tikzpicture}\\[6pt]
    {\scriptsize fin au centre, épais aux deux bords}

  \end{columns}

  \vspace{0.35cm}
  \centering\footnotesize
  Les trois profils sont \textbf{fins au centre} (point \textcolor{danger}{rouge},
  zone tactile) et \textbf{épais sur tout le pourtour}. L'exponentiel garde
  juste un fond central plus large et des bords qui montent plus vite.
  \emph{Épaisseur exagérée.}
\end{frame}



% ===================== Slide 14 : pourquoi epaissir les bords =====================
\begin{frame}{Pourquoi épaissir surtout les \emph{bords} ?}
  \begin{columns}[c]
  \column{0.5\textwidth}
    \centering
    \begin{tikzpicture}[scale=0.95]
      % impact au centre : grande portee
      \fill[gristik](0,2.2)--(-0.25,1.9)--(0.25,1.9)--cycle;
      \fill[gristik](4,2.2)--(3.75,1.9)--(4.25,1.9)--cycle;
      \draw[verre,very thick](0,2.2) .. controls (2,1.7) .. (4,2.2);
      \draw[-{Latex},danger](2,3)--(2,2.1);
      \node at (2,3.25){\scriptsize impact \textbf{centre} : grande portée $\Rightarrow$ souple};
      % impact au coin : petite portee
      \fill[gristik](0,0.4)--(-0.2,0.15)--(0.2,0.15)--cycle;
      \fill[gristik](1.4,0.4)--(1.2,0.15)--(1.6,0.15)--cycle;
      \draw[verre,very thick](0,0.4) .. controls (0.7,0.0) .. (1.4,0.4);
      \draw[-{Latex},danger](0.7,1.1)--(0.7,0.35);
      \node[align=center] at (2.7,0.55){\scriptsize impact \textbf{coin} : petite portée\\\scriptsize $\Rightarrow$ très contraint};
    \end{tikzpicture}
  \column{0.5\textwidth}
    \begin{itemize}
      \item Au \textbf{coin}, l'écran est tenu sur une petite portée
            $\Rightarrow$ la contrainte y est la plus forte.
      \item C'est exactement là que les fissures démarrent en pratique.
      \item Donc : \textcolor{verre}{\textbf{mettre la matière aux bords}}
            et garder le centre fin pour le toucher.
    \end{itemize}
  \end{columns}
\end{frame}

% ===================== Slide 15 : principe Monte-Carlo (bullets) =====================
\begin{frame}{Étape « simulation » --- la méthode de Monte-Carlo}
  \textbf{Idée :} chaque chute est différente (hauteur, angle, point d'impact).
  On en simule donc un \textbf{grand nombre au hasard} et on compte la proportion
  qui casse.
  \vspace{0.35cm}
  \begin{itemize}
    \item[\textcolor{verre}{\textbf{1.}}] \textbf{Tirer} une chute au hasard :
          hauteur $h$ et angle d'impact (selon une loi réaliste).
    \item[\textcolor{verre}{\textbf{2.}}] \textbf{Calculer} la contrainte maximale
          $\sigma_{\max}$ avec le modèle « poutre ».
    \item[\textcolor{verre}{\textbf{3.}}] \textbf{Comparer} au seuil de rupture :
          la dalle casse si $\sigma_{\max} > \sigma_c$.
    \item[\textcolor{verre}{\textbf{4.}}] \textbf{Répéter} l'opération $N=5000$ fois.
    \item[\textcolor{verre}{\textbf{5.}}] \textbf{Estimer} la probabilité de rupture :
      \[ P(\text{rupture}) \;\approx\; \frac{\text{nombre de casses}}{N} \]
  \end{itemize}
  \vspace{0.15cm}
  \begin{block}{Pourquoi ça marche}
  \footnotesize Plus $N$ est grand, plus cette proportion s'approche de la vraie
  probabilité : c'est la \textbf{loi des grands nombres}.
  \end{block}
\end{frame}

% ===================== Slide 16a : angle d'impact (inclinaison) 
\begin{frame}{Quel angle d'impact ? (inclinaison par rapport au sol)}
  \footnotesize
  L'angle $\theta$ est l'inclinaison de l'écran par rapport au sol au moment du choc.
  \vspace{0.5cm}
  \begin{center}
  \begin{minipage}{0.3\textwidth}\centering
    \begin{tikzpicture}[scale=0.8]
      \draw[thick](-0.7,0)--(1.9,0);
      \foreach \x in {-0.6,-0.3,0,0.3,0.6,0.9,1.2,1.5,1.8}{\draw[gristik](\x,0)--(\x-0.18,-0.18);}
      \draw[verre,thick,fill=verre!12,rounded corners=1pt](0,0) rectangle (1.5,0.16);
    \end{tikzpicture}\\[4pt]
    \scriptsize $\theta\approx 0^\circ$ (à plat)
  \end{minipage}%
  \hfill
  \begin{minipage}{0.33\textwidth}\centering
    \begin{tikzpicture}[scale=0.8]
      \draw[thick](-0.7,0)--(1.9,0);
      \foreach \x in {-0.6,-0.3,0,0.3,0.6,0.9,1.2,1.5,1.8}{\draw[gristik](\x,0)--(\x-0.18,-0.18);}
      \begin{scope}[rotate=45]
        \draw[verre,thick,fill=verre!12,rounded corners=1pt](0,0) rectangle (1.5,0.16);
      \end{scope}
      \draw[danger](0.55,0) arc (0:45:0.55) node[above,font=\tiny,danger]{$45^\circ$};
    \end{tikzpicture}\\[4pt]
    \scriptsize $\theta\approx 45^\circ$
  \end{minipage}%
  \hfill
  \begin{minipage}{0.33\textwidth}\centering
    \begin{tikzpicture}[scale=0.8]
      \draw[thick](-0.7,0)--(1.9,0);
      \foreach \x in {-0.6,-0.3,0,0.3,0.6,0.9,1.2,1.5,1.8}{\draw[gristik](\x,0)--(\x-0.18,-0.18);}
      \begin{scope}[rotate=80]
        \draw[danger,thick,fill=danger!12,rounded corners=1pt](0,0) rectangle (1.5,0.16);
      \end{scope}
      \draw[danger](0.5,0) arc (0:80:0.5) node[above,font=\tiny,danger]{$80^\circ$};
    \end{tikzpicture}\\[4pt]
    \scriptsize \textcolor{danger}{$\theta\approx 80^\circ$ (coin)}
  \end{minipage}
  \end{center}
  \vspace{0.3cm}
  \par\centering\footnotesize
  $\theta$ petit : choc à plat (peu dangereux)\;$\rightarrow$\;$\theta$ grand : choc sur le coin (dangereux).
\end{frame}


% ===================== Slide 16b : loi de l'angle (gaussienne) =====================
\begin{frame}{On tire $\theta$ selon une loi réaliste}
  \footnotesize
  \begin{itemize}
    \item Tirage d'une \textbf{gaussienne} $\mu=65^\circ$, $\sigma=15^\circ$, tronquée sur $[10^\circ,90^\circ]$.
    \item Les grands angles dominent : ce sont les chutes dangereuses.
  \end{itemize}
  \vspace{0.2cm}
  \centering
  \begin{tikzpicture}[scale=1]
    \draw[-{Latex}](0,0)--(9,0) node[right,font=\scriptsize]{$\theta\,(^\circ)$};
    \draw[-{Latex}](0,0)--(0,3.7) node[above,font=\scriptsize]{densité};
    \draw[verre,very thick,domain=10:90,samples=120,smooth]
      plot ({(\x-10)*8/80},{3*exp(-(\x-65)*(\x-65)/450)});
    % moyenne : libelle attache au haut du trait
    \draw[danger,dashed](5.5,0)--(5.5,3.0) node[above,font=\scriptsize,danger]{$\mu=65^\circ$};
    % ecart-type : libelle attache a la fleche
    \draw[ok,dashed](7,0)--(7,1.82);
    \draw[{Latex}-{Latex},ok,thick](5.5,1.25)--(7,1.25) node[midway,above,font=\scriptsize,ok]{$\sigma=15^\circ$};
    % troncature : libelle attache a un point (forme \draw ... node)
    \draw (1.0,3.0) node[gristik,font=\scriptsize,right]{tronquée sur $[10^\circ,90^\circ]$};
    % graduations : libelles attaches aux ticks
    \foreach \x/\lab in {0/10,4/50,5.5/65,7/80,8/90}{
      \draw (\x,0)--(\x,-0.12) node[below,font=\tiny]{\lab};}
  \end{tikzpicture}
\end{frame}




% ===================== Slide 17 : resultats =====================
\begin{frame}{Résultats : le profil change tout}
  \begin{columns}[c]
  \column{0.55\textwidth}
    \begin{center}
    \begin{tikzpicture}[yscale=0.042,xscale=1.05]
      \draw[-{Latex}](0,0)--(0,70) node[above,font=\scriptsize]{$P$ (\%)};
      \draw(0,0)--(5.1,0);
      \foreach \y in {0,20,40,60}{
        \draw(0,\y)--(-0.12,\y) node[left,font=\tiny]{\y};
        \draw[gray!25](0,\y)--(5.1,\y);}
      \fill[danger!70](0.4,0) rectangle (1.4,59);
      \fill[ok!70](2.0,0) rectangle (3.0,13);
      \fill[verre!70](3.6,0) rectangle (4.6,10);
      
ode[font=\tiny] at (0.9,64){59\%};
      
ode[font=\tiny] at (2.5,18){13\%};
      
ode[font=\tiny] at (4.1,15){10\%};
      
ode[font=\tiny,below] at (0.9,0){Uniforme};
      
ode[font=\tiny,below] at (2.5,0){Parabol.};
      
ode[font=\tiny,below] at (4.1,0){Expon.};
      
ode[font=\scriptsize] at (2.5,78){Risque de casse par profil};
    \end{tikzpicture}
    \end{center}
  \column{0.45\textwidth}
    \small
    \begin{tabular}{lc}
    \hline
    Profil & $P(\text{rupture})$ \\
    \hline
    Uniforme $0{,}5$~mm & $\approx 59\%$ \\
    Parabolique ($0{,}5\!\to\!1{,}2$) & $\approx 13\%$ \\
    Exponentiel & $\approx 10\%$ \\
    \hline
    \end{tabular}
    \vspace{0.3cm}
    \begin{beamercolorbox}[sep=5pt,rounded=true]{block body}
    \centering\footnotesize À centre identique, épaissir les bords
    \textbf{divise par 4 à 5} le risque de casse.
    \end{beamercolorbox}
    \footnotesize\textit{(valeurs de simulation, $N=5000$, ordres de grandeur)}
  \end{columns}
\end{frame}

% ===================== Slide 18 : optimisation =====================
\begin{frame}{Étape « optimisation » --- le meilleur compromis}
  \begin{columns}[c]
  \column{0.55\textwidth}
    \centering
    \begin{tikzpicture}[scale=0.62]
      % --- carte de chaleur P(e_c, e_b) dessinee en TikZ ---
      \foreach \i in {0,...,5}{
        \foreach \j in {0,...,3}{
          \pgfmathtruncatemacro\P{55 - \i*9 - \j*2}
          \pgfmathtruncatemacro\g{max(0,min(100,\P*1.6))}
          \fill[red!\g!green!88] (\i,\j) rectangle ++(1,1);
          
ode[font=\tiny,white] at (\i+0.5,\j+0.5){\P};
        }
      }
      \draw[step=1,gray!60,thin](0,0) grid (6,4);
      % zone interdite par la masse (coin haut-droit)
      \fill[pattern=north east lines,pattern color=black](5,3) rectangle (6,4);
      
ode[font=\tiny,align=center] at (5.5,3.5){};
      % optimum
      
ode[star,star points=5,star point ratio=2.2,fill=yellow,draw=black,
            scale=0.7] (opt) at (5.5,2.5){};
      
ode[font=\tiny,below=1pt of opt]{optimum};
      % axes
      \foreach \i/\lab in {0/0.5,1/0.8,2/1.1,3/1.4,4/1.7,5/2.0}{
        
ode[font=\tiny,below] at (\i+0.5,0){\lab};}
      
ode[font=\scriptsize,below] at (3,-0.55){$e_b$ (bord, mm)};
      \foreach \j/\lab in {0/0.2,1/0.3,2/0.4,3/0.5}{
        
ode[font=\tiny,left] at (0,\j+0.5){\lab};}
      
ode[font=\scriptsize,rotate=90] at (-1.0,2){$e_c$ (centre, mm)};
      
ode[font=\tiny,align=center] at (3,4.4)
        {\textcolor{danger}{rouge = risque fort}\ \ \textcolor{ok}{vert = risque faible}\ \ (valeurs en \%)};
    \end{tikzpicture}
  \column{0.45\textwidth}
    \small On cherche $(e_c,e_b)$ qui minimise $P$ sous deux contraintes :
    \begin{itemize}\footnotesize
      \item \textbf{toucher} : $e_c\le 0{,}5$~mm
      \item \textbf{masse} : $\le 2{,}5\times$ l'uniforme
    \end{itemize}
    \vspace{0.2cm}
    \begin{block}{Optimum trouvé}
    \centering\footnotesize centre fin + bords épais ($e_b\approx 2$~mm)\\
    $\Rightarrow P(\text{rupture})\approx 3\%$
    \end{block}
    \footnotesize Les deux contraintes sont \textbf{saturées} : c'est un vrai
    compromis « sécurité contre masse ».
  \end{columns}
\end{frame}


% =================================================================
%  SEQUENCE : Validation experimentale par CIRCUIT RLC
%  (verifie sigma_max ~ 1/sqrt(e) avec des grandeurs electriques)
% =================================================================

% ---------- 1) L'idee ----------
\begin{frame}{Validation : remplacer la mécanique par un circuit}
  \begin{columns}[c]
  \column{0.55\textwidth}
    \footnotesize
    \begin{itemize}
      \item Mesurer $\sigma$ et $e$ sur du verre est difficile.
      \item Mais l'oscillateur mécanique et un circuit \textbf{RLC série} suivent
            la \textbf{même équation}.
      \item On reproduit donc le choc avec une \textbf{impulsion électrique} et on
            vérifie la loi $\sigma_{\max}\propto 1/\sqrt{e}$ \textbf{uniquement avec
            des grandeurs électriques}.
    \end{itemize}
  \column{0.45\textwidth}
    \centering
    \begin{tikzpicture}[scale=0.9]
      \draw[decorate,decoration={coil,segment length=3pt,amplitude=3pt},verre,thick](0,1.2)--(1.2,1.2);
      \draw[fill=verre!12,thick](1.2,0.9) rectangle (1.9,1.5);
      
ode[font=\scriptsize] at (0.95,0.7){ressort + masse};
      \draw[-{Latex},thick](2.2,1.2)--(3.0,1.2);
      
ode[font=\scriptsize] at (2.6,1.5){$\equiv$};
      \draw[ok,very thick](3.3,1.35)--(3.9,1.35) (3.3,1.05)--(3.9,1.05);
      
ode[font=\scriptsize] at (3.6,0.7){circuit RLC};
    \end{tikzpicture}
  \end{columns}
\end{frame}

% ---------- 2) Le dictionnaire de l'analogie ----------
\begin{frame}{Le « dictionnaire » de l'analogie}
  \begin{columns}[c]
  \column{0.5\textwidth}
    \centering\footnotesize
    Même équation du 2\textsuperscript{nd} ordre :
    \[ m\ddot{x}+c\dot{x}+kx=F(t) \]
    \[ L\ddot{q}+R\dot{q}+\tfrac{1}{C}q=e(t) \]
    \vspace{0.2cm}
    \begin{tabular}{ll}
    \hline
    Mécanique & Électrique \\
    \hline
    masse $m$ & inductance $L$ \\
    amortissement $c$ & résistance $R$ \\
    raideur $k$ & $1/C$ \\
    fléche $x$ & charge $q$ \\
    vitesse $\dot x$ & courant $i$ \\
    choc $F$ & impulsion de tension \\
    \hline
    \end{tabular}
  \column{0.5\textwidth}
    \centering
    \begin{tikzpicture}[scale=0.85]
      % masse-ressort-amortisseur
      \draw[thick](0,2.4)--(2.4,2.4);
      \foreach \xx in {0,0.3,...,2.4}{\draw[gristik](\xx,2.4)--(\xx-0.2,2.6);}
      \draw[decorate,decoration={coil,segment length=3pt,amplitude=4pt},verre,thick](0.4,2.4)--(0.4,1.4);
      \draw[danger,thick](1.0,2.4)--(1.0,1.9) (0.85,1.9) rectangle (1.15,1.5);
      \draw[fill=verre!12,thick](0.1,0.8) rectangle (1.3,1.4);
      
ode[font=\scriptsize] at (0.7,1.1){$m$};
      
ode[verre,font=\scriptsize] at (0.0,1.9){$k$};
      
ode[danger,font=\scriptsize] at (1.45,1.7){$c$};
    \end{tikzpicture}
  \end{columns}
\end{frame}

\begin{frame}{Le montage : RLC série + générateur d'impulsions}
  \centering
  \begin{tikzpicture}[scale=1.2, every node/.style={font=\scriptsize}]
    % --- fils de la boucle (avec coupures pour les composants) ---
    \draw[thick] (0,0.7) -- (0,0) -- (4,0);          % bas
    \draw[thick] (0,1.3) -- (0,2) -- (0.6,2);        % haut-gauche
    \draw[thick] (1.6,2) -- (2.2,2);                 % entre L et R
    \draw[thick] (3.0,2) -- (4,2) -- (4,1.18);       % vers la capacite
    \draw[thick] (4,0.92) -- (4,0);                  % capacite -> bas

    % --- generateur d'impulsions ---
    \draw[thick] (0,1) circle (0.3);
    \draw[thick] (-0.12,0.95) -- (-0.04,0.95) -- (-0.04,1.12)
                 -- (0.04,1.12) -- (0.04,0.9) -- (0.12,0.9);
    
ode[left] at (-0.45,1) {GBF};

    % --- inductance L ---
    \draw[verre,thick,decorate,decoration={coil,segment length=4pt,amplitude=4pt}]
         (0.6,2) -- (1.6,2);
    
ode[verre] at (1.1,2.45) {$L$};

    % --- resistance R ---
    \draw[danger,thick] (2.2,1.8) rectangle (3.0,2.2);
    
ode[danger] at (2.6,2.45) {$R$};

    % --- condensateur C ---
    \draw[ok,very thick] (3.65,1.18) -- (4.35,1.18);
    \draw[ok,very thick] (3.65,0.92) -- (4.35,0.92);
    
ode[ok,right] at (4.45,1.05) {$C$};

    % --- oscilloscope sur C ---
    \draw[gray,dashed] (4.35,1.18) -- (5.2,1.6);
    \draw[gray,dashed] (4.35,0.92) -- (5.2,0.5);
    \draw[thick,rounded corners=2pt] (5.2,0.45) rectangle (6.4,1.65);
    
ode at (5.8,1.05) {$u_C(t)$};
    
ode at (5.8,1.95) {oscilloscope};
  \end{tikzpicture}

  \vspace{0.3cm}
  \par\footnotesize On fixe $L$, $R$ et l'\textbf{impulsion} ; on lit $u_C(t)$ sur
  l'oscilloscope. Le condensateur $C$ représente l'épaisseur.
\end{frame}




% ---------- 5) Protocole + lecture oscillo ----------
\begin{frame}{Protocole de mesure}
  \begin{columns}[c]
  \column{0.5\textwidth}
    \footnotesize
    \begin{enumerate}
      \item Fixer $L$, $R$, et l'amplitude/durée de l'\textbf{impulsion}.
      \item Pour chaque capacité $C_i$ :
      \begin{itemize}\scriptsize
        \item envoyer l'impulsion,
        \item relever $u_{C,\max}$ (1\textsuperscript{er} pic) et la
              pseudo-période $T$.
      \end{itemize}
      \item Calculer $e_i\propto C_i^{-1/3}$, $\;q_{\max}=C_i\,u_{C,\max}$,
            $\;\sigma_i\propto e_i\,q_{\max}$.
      \item Tracer $\sigma$ en fonction de $e$ (échelle log-log).
    \end{enumerate}
  \column{0.5\textwidth}
    \centering {\scriptsize régime transitoire lu à l'oscilloscope}\\[2pt]
    \begin{tikzpicture}[scale=0.9]
      \draw[-{Latex}](0,0)--(3.6,0) node[right,font=\tiny]{$t$};
      \draw[-{Latex}](0,-1)--(0,1.4) node[above,font=\tiny]{$u_C$};
      \draw[verre,thick,smooth,domain=0:3.3,samples=200]
            plot(\x,{1.1*exp(-0.7*\x)*cos(\x*520)});
      \draw[danger,dashed](0,1.1)--(0.35,1.1);
      
ode[danger,font=\tiny,left] at (0,1.1){$u_{C,\max}$};
      \draw[{Latex}-{Latex}](0.35,-0.75)--(1.05,-0.75) node[midway,below,font=\tiny]{$T$};
    \end{tikzpicture}
  \end{columns}
\end{frame}

% ---------- 6) Resultat attendu ----------
\begin{frame}{Résultat attendu : la pente $-1/2$}
  \begin{columns}[c]
  \column{0.5\textwidth}
    \centering
    \begin{tikzpicture}[scale=0.95]
      \draw[-{Latex}](0,0)--(4,0) node[right,font=\scriptsize]{$\ln e$};
      \draw[-{Latex}](0,0)--(0,3) node[above,font=\scriptsize]{$\ln \sigma$};
      \draw[verre,very thick](0.4,2.5)--(3.6,0.9);
      \foreach \x/\y in {0.6/2.4,1.2/2.1,1.9/1.75,2.6/1.4,3.3/1.05}{
        \fill[danger](\x,\y) circle (2pt);}
      
ode[verre,font=\scriptsize] at (2.6,2.3){pente $\approx -\tfrac12$};
    \end{tikzpicture}
  \column{0.5\textwidth}
    \footnotesize
    \begin{block}{Conclusion de l'expérience}
    Les points $(\ln e,\ln\sigma)$ s'alignent sur une droite de pente
    $\boxed{-\tfrac12}$ : la loi $\sigma_{\max}\propto e^{-1/2}$ est
    \textbf{vérifiée}, uniquement avec des mesures électriques.
    \end{block}
    \scriptsize C'est la dynamique du choc (oscillateur amorti) qui est validée,
    pas le verre lui-même.
  \end{columns}
\end{frame}

% ---------- 7) Tableau + limites ----------
\begin{frame}{Tableau de mesures et limites}
  \centering\scriptsize
  \begin{tabular}{cccccc}
  \hline
  $C$ (nF) & $e\propto C^{-1/3}$ (u.a.) & $u_{C,\max}$ (V) & $q_{\max}=C\,u_{C,\max}$ & $\sigma\propto e\,q_{\max}$ & $\ln\sigma$ \\
  \hline
  \dots & \dots & \dots & \dots & \dots & \dots \\
  \dots & \dots & \dots & \dots & \dots & \dots \\
  \dots & \dots & \dots & \dots & \dots & \dots \\
  \hline
  \end{tabular}
  \vspace{0.4cm}
  \begin{columns}[T]
  \column{0.5\textwidth}
    \footnotesize\textbf{Ce qui est validé}
    \begin{itemize}\scriptsize
      \item la dépendance $\sigma_{\max}\propto e^{-1/2}$,
      \item le régime transitoire d'un système du 2\textsuperscript{nd} ordre.
    \end{itemize}
  \column{0.5\textwidth}
    \footnotesize\textbf{Limites}
    \begin{itemize}\scriptsize
      \item l'analogie reproduit la \emph{dynamique}, pas la géométrie ni la
            rupture ;
      \item garder $R$ faible (régime pseudo-périodique) ;
      \item impulsion brève devant $T$ (approximation du choc).
    \end{itemize}
  \end{columns}
\end{frame}


% ===================== Slide 20 : limite 1D vs 2D =====================
\begin{frame}{Limite 1 --- poutre (1D) ou plaque (2D) ?}
  \begin{columns}[c]
  \column{0.5\textwidth}
    \centering
    \begin{tikzpicture}[scale=0.9]
      % poutre 1D
      \draw[verre,very thick](0,2.6) .. controls (1.3,1.9) .. (2.6,2.6);
      \draw(0,2.75)--(0,2.45)(2.6,2.75)--(2.6,2.45);
      \node at (1.3,3.1){\scriptsize Poutre 1D : courbure 1 direction};
      % plaque 2D : ondes concentriques
      \draw(0,-0.2) rectangle (2.6,1.4);
      \foreach \rr in {0.25,0.5,0.75,1.0}{\draw[verre](1.3,0.6) circle (\rr);}
      \fill(1.3,0.6) circle (1pt);
      \node at (1.3,-0.6){\scriptsize Plaque 2D : ondes dans 2 directions};
    \end{tikzpicture}
  \column{0.5\textwidth}
    \footnotesize
    \begin{itemize}
      \item Les particules bougent \textbf{perpendiculairement} à l'écran
            (flexion), pas « en 2D ».
      \item « 2D » = la déformation s'\textbf{étale sur la surface} (comme des
            ronds dans l'eau).
      \item La poutre ignore une direction $\Rightarrow$ elle
            \textbf{surestime} un peu $\sigma$ (prudent).
      \item Mais $k\propto e^3$ reste vrai $\Rightarrow$ \textbf{même
            conclusion}.
    \end{itemize}
  \end{columns}
\end{frame}

% ===================== Slide 21 : limite fabrication =====================
\begin{frame}{Limite 2 --- est-ce fabricable ?}
  \begin{columns}[c]
  \column{0.5\textwidth}
    \centering
    \begin{tikzpicture}[scale=1]
      % gradient continu (ideal)
      \draw[verre,thick](0,0)--(3,0);
      \draw[verre,thick](0,0.7) .. controls (1.5,0.1) .. (3,0.7);
      \draw[verre](0,0)--(0,0.7)(3,0)--(3,0.7);
      \node at (1.5,1.1){\scriptsize idéal : gradient continu};
      % paliers (realiste)
      \begin{scope}[yshift=-1.7cm]
        \draw[ok,thick](0,0)--(3,0);
        \draw[ok,thick](0,0.6)--(0.8,0.6)--(0.8,0.3)--(2.2,0.3)
              --(2.2,0.6)--(3,0.6);
        \draw[ok](0,0)--(0,0.6)(3,0)--(3,0.6);
        \node at (1.5,1.0){\scriptsize réaliste : paliers (marches)};
      \end{scope}
    \end{tikzpicture}
  \column{0.5\textwidth}
    \footnotesize
    \begin{itemize}
      \item Un gradient \textbf{continu} est difficile à produire (trempe en
            bain, usinage des bords).
      \item Alternatives crédibles :
      \begin{itemize}\scriptsize
        \item profil en \textbf{paliers} (marches discrètes) ;
        \item \textbf{renforcer la trempe} aux bords plutôt qu'ajouter de la
              matière (masse constante).
      \end{itemize}
      \item Le seuil $\sigma_c$ est \textbf{statistique} : notre $P$ donne un
            ordre de grandeur, pas une garantie.
    \end{itemize}
  \end{columns}
\end{frame}

% ===================== Slide 22 : conclusion =====================
\begin{frame}{Conclusion}
  \begin{columns}[c]
  \column{0.5\textwidth}
    \begin{block}{Réponse à la question}
    \footnotesize Le profil optimal est \textbf{fin au centre, épais aux
    bords}. Il fait passer le risque de casse d'environ
    \textcolor{danger}{$59\%$} (écran uniforme fin) à environ
    \textcolor{ok}{$3$--$13\%$}, à toucher préservé.
    \end{block}
    \footnotesize Tout repose sur une relation simple :
    \[ \sigma_{\max}\propto \frac{1}{\sqrt{e}} \]
    obtenue avec \textbf{uniquement} énergie + ressort + flexion d'une poutre.
  \column{0.5\textwidth}
    \centering
    \begin{tikzpicture}[scale=0.95]
      % deux barres P
      \draw[fill=danger!30](0,0) rectangle (0.9,3.0);
      \node[align=center] at (0.45,-0.45){\scriptsize uniforme\\\scriptsize $\approx 59\%$};
      \draw[fill=ok!35](1.8,0) rectangle (2.7,0.65);
      \node[align=center] at (2.25,-0.45){\scriptsize optimisé\\\scriptsize $\approx 3$--$13\%$};
      \draw[-{Latex},thick](1.0,2.4) .. controls (1.4,2.4) .. (1.7,1.0);
      \node at (1.45,2.7){\scriptsize $\div 4$ à $5$};
    \end{tikzpicture}
  \end{columns}
\end{frame}

% ===================== Slide 23 : ouvertures =====================
\begin{frame}{Ouvertures et bibliographie}
  \begin{columns}[T]
  \column{0.5\textwidth}
    \textbf{ }
    \footnotesize
    \begin{itemize}
      \item Modèle \textbf{plaque 2D} (Kirchhoff--Love) + contact de Hertz.
      \item Loi de \textbf{Weibull} pour la rupture (vraie statistique des
            défauts).
      \item Matériaux à \textbf{gradient de propriétés} (FGM) ;
            trempe différentielle aux bords.
    \end{itemize}
  \column{0.5\textwidth}
    \textbf{Sources}
    \scriptsize
    \begin{itemize}
      \item Timoshenko, \textit{Theory of Plates and Shells}, 1959.
      \item K.~L.~Johnson, \textit{Contact Mechanics}, 1985.
      \item W.~Weibull, \textit{J. Appl. Mech.}, 1951.
      \item R.~Gy, \textit{Mat. Sci. Eng. B}, 2008 (trempe par échange ionique).
      \item Corning Gorilla Glass --- fiches techniques.
    \end{itemize}
  \end{columns}
  \vspace{0.6cm}
  \begin{center}\Large\textcolor{verre}{Merci de votre attention.}\end{center}
\end{frame}

\begin{frame}[fragile,allowframebreaks]{Annexe --- Code Python}
\tiny
\begin{verbatim}
# -*- coding: utf-8 -*-
# TIPE - Modele SIMPLIFIE (poutre 1D)
# Probabilite de rupture par Monte-Carlo.

import numpy as np
import matplotlib.pyplot as plt

rng = np.random.default_rng(2026)

# --- constantes (verre trempe + smartphone) ---
E, G = 73e9, 9.81          # module d'Young (Pa), pesanteur
M_PHONE = 0.18             # masse du telephone (kg)
ETA = 0.05                 # fraction d'energie stockee en flexion
B = 0.020                  # largeur de la bande de verre (m)
SIG_MU, SIG_SD = 850e6, 120e6   # seuil de rupture gaussien (Pa)
KAPPA = 0.7                # surcontrainte au coin
L_C, L_B = 0.050, 0.025    # portee effective : centre / coin (m)
E0 = 0.5e-3                # epaisseur uniforme de reference (m)
N_MC = 5000

# --- profils d'epaisseur e(r), r de 0 (centre) a 1 (bord) ---
def uniforme(r, e0=E0):
    return np.full_like(np.asarray(r, float), e0)

def parabolique(r, e_c=0.5e-3, e_b=1.2e-3):
    return e_c + (e_b - e_c) * np.asarray(r, float) ** 2

def exponentiel(r, e_c=0.5e-3, e_b=1.2e-3):
    return e_c * np.exp(np.log(e_b / e_c) * np.asarray(r, float))

def epaisseur_moyenne(prof, **kw):
    r = np.linspace(0, 1, 400)
    return np.trapz(prof(r, **kw) * 2 * r, r) / np.trapz(2 * r, r)

# --- coeur du modele : contrainte de flexion (poutre 1D) ---
def contrainte(e, L, Ek, Kc):
    I = B * e ** 3 / 12.0           # moment quadratique (RDM)
    k = 48.0 * E * I / L ** 3       # raideur poutre : k ~ e^3
    x = np.sqrt(2.0 * ETA * Ek / k) # 1/2 k x^2 = eta Ek -> fleche max
    return 6.0 * E * e * x / L ** 2 * Kc

def gauss_tronquee(mu, s, lo, hi, n):
    out = np.empty(n); k = 0
    while k < n:
        x = rng.normal(mu, s, n - k); x = x[(x >= lo) & (x <= hi)]
        out[k:k + x.size] = x; k += x.size
    return out

def monte_carlo(prof, N=N_MC, **kw):
    theta = gauss_tronquee(65, 15, 10, 90, N)     # angle de chute (deg)
    h = gauss_tronquee(1.0, 0.25, 0.30, 2.0, N)   # hauteur (m)
    v = np.sqrt(2 * G * h); Ek = 0.5 * M_PHONE * v ** 2
    r = np.clip(theta / 90 + rng.normal(0, 0.05, N), 0, 1)
    L = L_C - (L_C - L_B) * r                     # portee selon la zone
    Kc = 1 + KAPPA * r                            # coin = surcontrainte
    e = prof(r, **kw)
    sigma = contrainte(e, L, Ek, Kc)
    sigma_c = np.maximum(rng.normal(SIG_MU, SIG_SD, N), 1e6)
    return (sigma > sigma_c).mean(), sigma

def optimise():
    emu = epaisseur_moyenne(uniforme)
    ec_g = np.linspace(0.2e-3, 0.5e-3, 7)
    eb_g = np.linspace(0.5e-3, 2.0e-3, 16)
    P = np.full((ec_g.size, eb_g.size), np.nan); best = None
    for i, ec in enumerate(ec_g):
        for j, eb in enumerate(eb_g):
            if eb < ec or epaisseur_moyenne(parabolique, e_c=ec, e_b=eb) > 2.5 * emu:
                continue
            p, _ = monte_carlo(parabolique, N=3000, e_c=ec, e_b=eb)
            P[i, j] = p
            if best is None or p < best[0]:
                best = (p, ec, eb)
    return ec_g, eb_g, P, best

if __name__ == "__main__":
    v1 = np.sqrt(2 * G * 1.0); Ek1 = 0.5 * M_PHONE * v1 ** 2
    print("v=%.2f m/s  Ek=%.2f J" % (v1, Ek1))
    for nom, f, kw in [("Uniforme", uniforme, {"e0": E0}),
                       ("Parabolique", parabolique, {"e_c":0.5e-3,"e_b":1.2e-3}),
                       ("Exponentiel", exponentiel, {"e_c":0.5e-3,"e_b":1.2e-3})]:
        p, s = monte_carlo(f, **kw)
        print(nom, "P=%.1f%%" % (p*100))
    ec_g, eb_g, P, best = optimise()
    print("Optimum P=%.1f%% ec=%.2fmm eb=%.2fmm"
          % (best[0]*100, best[1]*1e3, best[2]*1e3))
\end{verbatim}
\end{frame}

\end{document}


